 \begin{center}  МИНОБРНАУКИ РОССИИ\\Федеральное государственное бюджетное образовательное\\
  учреждение высшего образования\\«Ижевский государственный технический университет \\имени М.Т. Калашникова»\\  Факультет <<Приборостроительный>>\\Кафедра «Радиотехника»\\
 \end{center} % вертикальные пробельные отступы, вычисляемые автоматически
 \vfill\vfill 
 \begin{center}{\bf\MakeTextUppercase{Лабораторная работа №1}}\\  тема: <<Изучение основ построения IP-сетей>>\\по дисциплине: <<Сетевые информационные технологии>>\\Вариант 2
 \end{center} \vfill
 \vfill \noindent
 \begin{tabular}{p{0.7\linewidth}l}  Выполнил & \\cтудент гр. М23-762-1 & Асланян Д.Б.\\
  & \\Проверил& \\  к.т.н., доцент & Ардашев Д.В.\\& \\
 \end{tabular} \vfill
 \begin{center}  Ижевск \the\year{}
 \end{center} \newpage
 \pagestyle{plain} \parindent=1.25cm
